\documentclass[11pt]{article}
\usepackage{amsmath,amssymb,amsthm,hyperref}
\newtheorem{theorem}{Theorem}
\newtheorem{lemma}[theorem]{Lemma}
\title{A deterministic lower bound for the binomial triple projection}
\author{Research note; independently reviewed by two other agents}
\date{30 September 2026}
\begin{document}
\maketitle

Let $\rho_1,\ldots,\rho_m$ be arbitrary permutations of $[n]$, with no
randomness or independence assumption. On distinct labels put
\[
 \delta_i(a,b,c)=\frac12-\frac32
    \mathbf1\{\rho_i(b)<\rho_i(a),\rho_i(b)<\rho_i(c)\}.
\]
Write $d=n-3$, $p_i=(i-1)/(m-1)$, and
\[
 w_i(a)=\binom da p_i^a(1-p_i)^{d-a},\qquad
 h=\binom n3m^{-3}(1-1/m)^{n-3}.
\]
The triple contribution in the palindrome expansion is the following
function on target permutations:
\[
 S(\sigma)=h\sum_{a=0}^{d}\sum_{i=1}^m
 w_i(a)\delta_i(\sigma_{a+1},\sigma_{a+2},\sigma_{a+3}).
\]
The terminology ``projection'' refers to the decomposition for the
independent random model. The displayed function itself makes sense for
every deterministic block sequence.

\begin{theorem}\label{main}
There are absolute constants $c_0,c_1>0$ and $n_0$ such that, for
$n\ge n_0$ and
\[
 n\le m\le c_0 n^{3/2},
\]
every choice of the block permutations satisfies
\[
 \max_{\sigma\in S_n}S(\sigma)-\min_{\sigma\in S_n}S(\sigma)
 \ge c_1\frac{n^{13/4}}{m^{5/2}}.
\]
The range is witnessed by independent middle-pair swaps in disjoint
six-label windows in the central third of the target permutation.
\end{theorem}

\paragraph{Scope.}
This theorem concerns the displayed linear contribution. It is not a
lower bound for the complete ordering error of arbitrary deterministic
palindrome sequences. The other terms can depend strongly on the block
sequence; the probabilistic remainder bounds for independent blocks
cannot be used for a deterministic sequence.

\begin{lemma}[Central binomial weights]\label{weights}
For $d$ sufficiently large, $d/4\le a\le3d/4$, and $m\ge d+3$, put
$Z_a=\sum_iw_i(a)$. There are absolute constants $C,c>0$ such that
\[
 c\frac md\le Z_a\le C\frac md,
 \qquad c\frac m{d^{3/2}}\le\sum_iw_i(a)^2
                           \le C\frac m{d^{3/2}},
\]
\[
 \qquad
 \sum_i|w_i(a+r)-w_i(a)|\le C\frac m{d^{3/2}}
 \quad (|r|\le3),
\]
provided the shifted indices also lie in $[d/4,3d/4]$. In particular,
the normalized weights $\mu_i=w_i(a)/Z_a$ satisfy
$\sum_i\mu_i^2\asymp\sqrt d/m$.
\end{lemma}
\begin{proof}
Let $f_a(p)=\binom da p^a(1-p)^{d-a}$. It is unimodal, has integral
$1/(d+1)$, and, by Stirling's bounds, maximum at most $C/\sqrt d$
uniformly over central $a$. For a function $f$ of bounded variation,
its sum on the grid of mesh $1/(m-1)$ differs from $(m-1)\int f$ by
at most its total variation plus endpoint terms. Consequently the
error for $f_a$ is $O(d^{-1/2})$, proving the estimates for $Z_a$.
For the squared weights, the beta integral and the same Stirling bounds
give
\[
 \int_0^1 f_a(p)^2\,dp
 =\frac{\binom da^2}{(2d+1)\binom{2d}{2a}}
 \asymp d^{-3/2}.
\]
The squared function is unimodal with maximum $O(d^{-1})$, so its
grid error is $O(d^{-1})$, smaller than its main term for $m\ge d+3$
and sufficiently large $d$. This proves the squared-weight estimate.

For neighboring indices, let
$g(p)=\binom{d+1}{a+1}p^{a+1}(1-p)^{d-a}$. Direct differentiation gives
\[
 g'(p)=(d+1)(f_a(p)-f_{a+1}(p)).
\]
Thus $\int|f_a-f_{a+1}|=2\max g/(d+1)=O(d^{-3/2})$.
The total variation of $|f_a-f_{a+1}|$ is at most the sum of the
variations of $f_a,f_{a+1}$, hence is $O(d^{-1/2})$. The grid estimate
therefore gives $O(md^{-3/2}+d^{-1/2})=O(md^{-3/2})$.
Telescoping proves the claim for all the stated shifts.
\end{proof}

\begin{proof}[Proof of Theorem~\ref{main}]
Use the weighted six-window theorem in the companion note
\emph{Weighted collision bounds for triple minima and six-window
swaps}: for a weighted permutation family on a set of $N$ labels with
$\kappa=\sum_i\mu_i^2$ and $N\ge C_*/\kappa$, some six distinct
labels have a middle-pair swap changing the sum of their four consecutive
triple scores by at least $c_*\sqrt\kappa$. Its constants are absolute,
and its conclusion is preserved on restriction to a label subset.

Fix $J=\lfloor n/40\rfloor$ disjoint six-position windows in the central
third of the target order; a harmless increase in $n_0$ ensures that
this is possible. Fill these windows successively with labels, always
using labels not used in earlier windows. At every step at least
$n-6J\ge n/2$ labels are available. If $a$ is the start of the first
of the four affected triples in the current window, all starts
$a,a+1,a+2,a+3$ are central.

Take the distribution $\mu_i=w_i(a)/Z_a$. By Lemma~\ref{weights},
$\kappa\asymp\sqrt n/m$. We can choose $c_0$ sufficiently small that
$n/2\ge C_*/\kappa$. Apply the weighted six-window theorem on the
unused labels. It supplies six labels and a swap with absolute
normalized score difference at least
$c n^{1/4}/\sqrt m$.

Replace each constant weight
$w_i(a)$ by the actual weight $w_i(a+r)$ in its corresponding triple,
for $r=0,1,2,3$. Lemma~\ref{weights} bounds the resulting normalized
change by
\[
 \frac2{Z_a}\sum_{r=0}^3\sum_i|w_i(a+r)-w_i(a)|
       \le \frac C{\sqrt n}.
\]
After decreasing $c_0$ once more, this error is less than half the
supplied swap difference. Thus the actual change in $S$ from this window
is at least
\[
 c\,h Z_a\frac{n^{1/4}}{\sqrt m}
 \ge c\,h\frac mn\frac{n^{1/4}}{\sqrt m}.
\]
Repeat for all windows and put the remaining labels arbitrarily in the
remaining positions. The two exterior labels on each side of each
swapped pair stay fixed, so all affected triples lie inside their
own window. The effects of the $J$ binary choices are exactly additive.
Hence the range is at least the sum of the $J$ absolute changes.
Using $J\ge cn$ and $h\ge c n^3/m^3$ proves
\[
 \operatorname{range}(S)
 \ge c n\frac{n^3}{m^3}\frac mn\frac{n^{1/4}}{\sqrt m}
 =c\frac{n^{13/4}}{m^{5/2}},
\]
as required.
\end{proof}
\end{document}
